\scanpage{4}
\mspage{3}
Si $F,F'$ sont deux figures, on dit que $F'$ est une \uline{sous-figure} de $F$, si \unsure{toute|tout} strate de $F'$ est une strate de $F$, i.e. $F'\subset F$. Donc \deleted{le support} on aura \deleted{\unsure{$\sigma(F)$}} \unsure{d|et} $\lvert F'\rvert\subset\lvert F\rvert$

\uline{NB} Si $X_i\in F'$, et $Y_j\in F$ est tel que $Y_j\subset X_i$, alors $Y_j\in F'$

\deleted{En effet, soit $y\in\deleted{\unsure{Y_j}}\operatorname{int}(Y_j)$, alors}

\deleted{La relation}

On écrit $F'\leq F$ pour cette relation, qui est une relation d'ordre dans $\mathcal{F}$.

L'ens des sous-figures de $F$ est en corr. 1-1 avec l'ens des parties \inserted{$\Phi$ ($=\Sigma_{F'}$)} de \deleted{\unsure{$F$|$\mathcal{F}$}} \inserted{$\Sigma_F$, (ens des strates)} qui ont la \unsure{propriét.|propriété} que \inserted{$i,j\in\Sigma_F$,} $i\in\deleted{\unsure{\Sigma|B}}\Phi$, $j\leq i\Rightarrow j\in\Phi$
(\deleted{idéaux} \inserted{``\unsure{fermées|fermés}''} \deleted{fermés} de l'ens \deleted{\unsure{ens}} ordonné, pour la top. \inserted{sur $\Sigma_F=\Sigma$} \unsure{ayant} comme base \deleted{de fermés} \deleted{d'ouverts \unsure{/} les} \inserted{d'ouverts les} $\Sigma_{\deleted{\unsure{>}}\unsure{\geq|\leq}i}$).

\deleted{\unsure{Ainsi|$F'\leq F$}, on a une relation d'ordre sur}
% slanted
\margin{\uline{11.6.} Je dirai plus tôt ``\uline{multistrate}''}

\uline{Figure élémentaire} : qui a une strate maximale i.e. $\Sigma_F$ a un plus grand élément. Soit \tikz[baseline=(el.base)]{\node[inner sep=1pt] (el) {$\mathcal{F}_{\text{él}}$};\draw (el.center) ellipse [x radius=1.05em,y radius=.78em];} l'ens des figures élémentaires. \deleted{Alors} Pour tte figure $F$, les sous-figures élémentaires sont en corr 1-1 avec les
